%%%%%%%%%%%%%%%%%%%%%%%%%%%%%%%%%%%%%%%%%%%%%%%%%
\section{The WSU Signal Chain}
\label{sec:wsu_signal_chain}
%%%%%%%%%%%%%%%%%%%%%%%%%%%%%%%%%%%%%%%%%%%%%%%%%

This section provides a high-level overview of the \ac{WSU} signal chain, from the Front-End receivers through the \ac{IF} signal conditioning and digitization stages, the high-speed data transmission system and the fiber optic transport infrastructure, to the digital signal processing performed in the \ac{ATAC} and \ac{TPGS}. 
It describes the upgraded system architecture that supports four wider bandwidth \ac{IF}  signals per antenna, enabled by the deployment of 4$\times$\ac{IF} bandwidth receivers, and covers signal conditioning, digitization, transmission, reference signal distribution, and site infrastructure integration.


%%%%%%%%%%%%%%%%%%%%%%%%%%%%%%%%%%%%%%%%%%%%%%%%%
\subsection{Signal Chain from the Front-Ends to the Correlator/Spectrometer}
\label{sec:WSUSignalChain}

%%%%%%%%%%%%%%%%%%%%%%%%%%%%%%%%%%%%%%%%%%%%%%%%%


This section describes the architecture of the \ac{WSU} signal chain, with a focus on the transition from analog to digital processing. Fig.~\ref{fig:wsu-system-block-diagram} presents the block diagram of both the legacy and \ac{WSU} signal chains, illustrating the concept of parallel deployment where both systems coexist to ensure a smooth transition. 

A general trend in modern high-performance receiver design is to digitize the analog signal as early as possible in the signal chain.
Signal processing is then performed in the digital domain whenever technically and economically feasible.
The digitizer is the boundary between the analog and digital domains.
Given the requirements of the ALMA \ac{WSU}, particularly in terms of (a) instantaneous bandwidth (the frequency range that can be observed simultaneously without retuning) and (b) signal quality, direct digitization of the RF signal remains challenging for \ac{ALMA} due to the extremely high sampling rate and dynamic range required up to 950~GHz. 

With significant advancements in digitizer technology, sampling and quantization at the first \ac{IF} signal stage has now become viable for \acs{ALMA}, even while meeting the objective of achieving four times the instantaneous bandwidth of the legacy system.
This technological breakthrough has simplified the analog signal chain, eliminating the need for a second \ac{IF} stage of the legacy \ac{BE} subsystem. Fig.~\ref{fig:antennaToCorrelator} introduces the near-real-time digital signal processing chain, which will be described in more detail in a later section.

The \ac{WSU} signal chain efficiently processes and delivers astronomical signals from the \ac{FE} receivers to the \ac{ATAC} and \ac{TPGS} subsystems.
The key objective is to digitize signals as early as possible, maintaining high fidelity while expanding the instantaneous bandwidth to four times that of the legacy system. \\

%%%%%%%%%%%%%%%%%%%%%
\begin{figure}[t!]
\centering
\includegraphics[width=\textwidth]{Figures/ALMA2030_Concept_of_Parallel_Deployment_2025May20.png}
\caption{WSU system block diagram and the concept of parallel deployment, with
the WSU signal path and components highlighted in green and the Legacy
signal path and components highlighted in blue. \textit{Note:} This figure presents a basic functional diagram of the \ac{WIFP}. In the current architecture, these functions are distributed between the IF
Switch and \ac{DGOT}, and the physical integration of \ac{DGOT} and \ac{DSPOT} is still under consideration. This figure should be regarded as provisional,
subject to future revision.
\label{fig:wsu-system-block-diagram}}
\end{figure}
%%%%%%%%%%%%%%%%%%%%%
% Notes for TOK
% Note: "ALMA WSU" instead of "ALMA 2030"?
%%%%%%%%%%%%%%%%%%%%%%%%%%%%%%%%%%%%%%%%%%%%%%%%%

%%%%%%%%%%%%%%%%%%%%%
\begin{figure}[hbt!]
\centering
\includegraphics[width=0.96\textwidth]{Figures/Antenna-to-Correlator-Diagram.png}
\caption{Simplified WSU system block diagram focusing on the near-real time digital signal path \cite{cosdd_c}.\label{fig:antennaToCorrelator}}
\end{figure}
%%%%%%%%%%%%%%%%%%%%%


\noindent
\textbf{Signal Reception and IF Down-conversion}

The process begins at the \ac{FE} subsystem, which houses up to ten receiver bands covering frequencies from 35~GHz to~950 GHz.
The \ac{IF} signals are generated by down-converting the received \ac{RF} signals with the \ac{LO1} signals, producing \ac{IF} signals spanning $2-20$~GHz, a significant expansion compared to the $4-12$~GHz range of the legacy system. 
The legacy receivers provide either four \ac{IF} outputs with at least 4~GHz bandwidth per sideband per polarization, or two outputs with 8~GHz bandwidth per polarization, depending on the receiver band.
In contrast, the wideband receivers implemented in the \ac{WSU} provide four \ac{IF} outputs, each with at least 8~GHz and up to 16~GHz bandwidth per sideband per polarization, depending on the receiver band.
Consequently, this represents a 4$\times$ increase in the processed bandwidth.


%%%%%%%%%%%%%%%%%%%%%%%%%%%%%%%%%%%%%%%%%%%%%%%%%
\noindent
\textbf{IF Signal Conditioning and Digitization}

The \ac{IF} signals pass through the \ac{WIFP}, which conditions, digitizes, and transmits them to the \ac{DTS}.
The IF Switch within the FE selects receiver bands, provides slope compensation for following cables, and adjustable gain (1dB steps). The IF signals are then connected to a bulkhead panel on the FE. From there they connect to the \ac{DGOT} unit. During parallel deployment a splitter will allow both the DGOT and the legacy IF Processors to receive the IF signals. The IF conditioner within the DGOT provides further amplification, gain adjustment and slope correction, and filtering $2-20$~GHz (the upper cut-off being particularly sharp as it provides anti-aliasing). 
For the legacy signal chain there is filtering within the legacy BE of $4-12$~GHz prior to downconversions to basebands of $2-4$~GHz which are provided to the legacy digitizers. The baseband anti-aliasing filters result in an effective instantaneous bandwidth of approximately 1.875~GHz per baseband per polarization, resulting in 3.75~GHz per sideband per polarization for \ac{2SB} receivers, and 7.5~GHz per polarization for \ac{SSB} and \ac{DSB} receivers.

For the \ac{WSU} path, the \ac{IF} signal chain is designed to support operation across an approximate $2-20$~GHz range, with specifications for the anti-aliasing filters defined in the \ac{WIFP} document: $F_{IFmin}$ $<$ 2~GHz, $F_{IFmax}$ $>$ 19.5~GHz, and $\le$1~dB insertion loss \cite{WIFPDescription}.
The \ac{DGOT} module then digitizes these signals at 6-bit resolution and 40~GSps per sideband per polarization.
In the \ac{WSU} signal chain, the digitizer resolution has been increased from 3~bits in the legacy system to 6~bits, significantly improving digitizer quantization efficiency and thereby contributing to the higher overall digital efficiency achieved across the \ac{WSU} signal chain \cite{orig_dig_ef_budget}.
The \ac{DGCTR} provides a 20~GHz clock derived from the 125~MHz reference, ensuring synchronization. 
Prior to digitization the IF signals are also split to drive a set of total power detectors covering a set of bandwidths within the IF.
The digitized \ac{IF} signals are subsequently processed in the \ac{DSPOT} module, where they are packetized and sent to four 400\,GbE optical transceivers provided by the \ac{DTS}.

%%%%%%%%%%%%%%%%%%%%%%%%%%%%%%%%%%%%%%%%%%%%%%%%%
\noindent
\textbf{High-Speed Data Transmission}

Once the signals are digitized, they are transmitted over the \ac{WSU} \ac{DTS} to the \ac{ATAC} and \ac{TPGS} subsystems.
The \ac{WSU} \ac{DTS} utilizes 400\,GbE optical links with \ac{DWDM} technology, multiplexing four independent data streams -- one per polarization sideband -- from each antenna. 
With a total raw data rate of 960~Gbps per antenna, the optical signals are amplified, transmitted over fiber distances of up to 80~km, and then demultiplexed for processing at \ac{ATAC} and \ac{TPGS}.

The fiber optic infrastructure that supports the \ac{WSU} \ac{DTS} consists of two segments: an existing fiber network between the antennas and the \ac{AOS} \ac{TB}, and a newly installed \ac{WSU} \ac{FOS} between the \ac{AOS} \ac{TB} and the \ac{OSF} (see Section~\ref{sec:FOS}).


%%%%%%%%%%%%%%%%%%%%%%%%%%%%%%%%%%%%%%%%%%%%%%%%%
\noindent
\textbf{Advanced Digital Signal Processing}

Using \ac{ATAC}, 7-m Array and 12-m Array data are processed in three stages using up to eight independent concurrent correlator subarrays. The ALMA Very Coarse Channelizers (\acs{AVCC}s) apply Oversampling Polyphase Filter Banks (\acs{OSPFB}s) to generate 222.222~MSps wide Frequency Slices (\acs{FS}s), each of which is then processed by an \ac{AFSP} which applies fringe tracking, delay compensation, and removal of LO Offset per antenna while resampling each \acs{FS} to produce 200.88~MHz of science quality bandwidth.
Finally, the \ac{CDP} combines adjacent Frequency Slices (\acs{FS}s) to form Science Spectral Windows, applies Van Vleck corrections, normalizes visibility data, and formats it for archival storage.

The \ac{TPGS} subsystem, built on Graphics Processing Units (\acs{GPU}s), is primarily designed to generate auto-correlation products for single dish observations from the four \ac{TP} antennas.
It utilizes an FFX-type signal processing architecture—comprising coarse channelization, fine channelization, and multiplication—aligned with the design of the \ac{ATAC} subsystem.
While its main function is auto-correlation processing, the \ac{TPGS} also performs cross-correlation calculations among the four antennas, used exclusively for pointing and focus calibrations.
Upon initial deployment at the \ac{OSF}, the \ac{TPGS} subsystem will support 4$\times$ \ac{IF} bandwidth (32~GHz per polarization) coverage.\\


%%%%%%%%%%%%%%%%%%%%%%%%%%%%%%%%%%%%%%%%%%%%%%%%%
\noindent
\textbf{Local Oscillator Subsystem and Time Synchronization}

To ensure phase coherence across all antennas, the \ac{WSU} signal chain operates with the legacy \ac{PLO} subsystem, managed by the \ac{CLOA}, \ac{LORR}, and \ac{LPR}.
Additionally, the \ac{DGCTR} module within the \ac{WIFP} ensures synchronization with an accuracy of $\le 25$~ps jitter, embedding precise time indices into digitized signals for correlation.

%%%%%%%%%%%%%%%%%%%%%%%%%%%%%%%%%%%%%%%%%%%%%%%%%
\noindent
\textbf{Delay Correction}

The \ac{ALMA} “Delay Server” software computes the delay model for the signal chain at the required cadence and distributes this information to the correlator subsystems and the antennas accordingly.
Fringe tracking and all of the delay compensation, including the cable delay of the \ac{LO} reference, is handled in the correlator subsystems (\ac{ATAC} and \ac{TPGS}) based on the received delay model and the absolute time contained in the timestamps of the digitized datastream.
Knowledge of the delay is needed at the antennas only to control the timing of the 180~degree Walsh modulation/demodulation so that the switching pattern is aligned when correlation is performed at the correlator.

%\newpage

%%%%%%%%%%%%%%%%%%%%%%%%%%%%%%%%%%%%%%%%%%%%%%%%%
\noindent
\textbf{LO Offset for Interference Mitigation and Sideband Suppression}

To mitigate internally generated interference and suppress the image sideband (especially relevant for legacy \ac{DSB} receivers, Bands 9 and 10), the \ac{WSU} signal employs an \ac{LO} offset technique.
This process is managed by the \ac{FLOOG}, which introduces a small frequency offset to the \ac{LO1} signal at each antenna.
While this offset does not eliminate noise from the unwanted sideband, it shifts correlated signals to a high fringe rate, effectively spreading them across a larger bandwidth and increasing noise incoherently.

In this approach, a frequency offset, $f_0$, is applied to \ac{LO1} by \ac{FLOOG} at each antenna and subsequently removed by \ac{ATAC} and \ac{TPGS}.
This ensures that the signal sideband remains at the correct frequency while the image sideband is shifted by $2f_0$ from its nominal value.
To prevent unwanted sideband signals from correlating across antennas, different values of $f_0$ are uniquely assigned to each antenna.
This configuration keeps all signal sidebands aligned in frequency while dispersing the image sidebands, reducing their coherence and minimizing their impact on observations.
As the \ac{LO} offset is removed through digital signal processing in the \ac{ATAC} and \ac{TPGS} subsystems, this approach requires that the high quality portion of the first channelization stage’s spectral profile is wide enough to accommodate all the necessary topocentric shifts of the \ac{FS}, including \ac{LO} offset removal and alignment of fine channels in adjacent \ac{FS} when producing wider spectral windows.

%%%%%%%%%%%%%%%%%%%%%%%%%%%%%%%%%%%%%%%%%%%%%%%%%
\noindent
\textbf{Walsh Phase Switching in the WSU Signal Chain}

The \ac{WSU} signal chain employs 180-degree Walsh phase switching to reject spurious signals before digitization.
The \ac{FLOOG} applies this phase switching based on orthogonal Walsh function patterns, with a cycle time of 16~ms.
Each antenna is assigned a unique Walsh pattern, which is later demodulated by the \ac{DSPOT} module in the \ac{WIFP}.
This ensures that desired signals correlate while unwanted signals are effectively canceled, mitigating spurious signals generated between the (first) down-converter and the \ac{DGOT} and suppressing digitizer DC offsets.

In contrast, 90-degree phase switching for sideband separation is not required in the \ac{WSU} signal chain to process the existing \ac{DSB} receivers (Bands 9 and 10), and thus it will not be supported. Sideband suppression in \ac{DSB} receivers will now employ the LO offset technique explained in the preceding subsection.


%%%%%%%%%%%%%%%%%%%%%%%%%%%%%%%%%%%%%%%%%%%%%%%%%
\subsection{Receivers}
\label{sec:receivers}

%%%%%%%%%%%%%%%%%%%%%%%%%%%%%%%%%%%%%%%%%%%%%%%%%

The ALMA \ac{FE} subsystem is designed to accommodate up to 10 receiver bands, covering most of the wavelength range from 8.5 to 0.32~mm ($35-950$~GHz).
Each band is designed to cover a tuning range that is approximately tailored to the atmospheric transmission windows.
Upon completion of ongoing development projects addressing the top-level scientific priorities of the \ac{WSU}, and as of Milestone~5, Bands 1, 2, 3, 4, 5, 6v2, 7v2, 8v2, 9, and 10 will be available. Their basic characteristics are outlined in Table~\ref{tab:rx_summary}.


%%%%%%%%%%%%%%%%%%%%%
% RID-4936/RID-4937/RID-5007
%%%%%%%%%%%%%%%%%%%%%
\begin{table}[h!]
    \centering
    \caption{ALMA Receiver Specifications}
    \label{tab:rx_summary}
    \begin{tabular}{llllllll}
        \hline
        Band & Frequency & Sideband & IF range & Inst. IF BW & $T_{rx}$ over 80\% & $T_{rx}$ at any\\
        & range (GHz) & mode & (GHz) & (GHz)$^4$ & of band (K)$^6$ & freq. (K)$^6$\\
        \hline
        1 & $35-50$ &  \acs{SSB} & $4-12$ & $8$ & $<28$ & $<32$\\
        2$^1$ & $67-116$ & \acs{2SB} & $2-18$ & $32$ & $<30^7$ & $<43$\\
        & & & & & $<41^8$ &\\
        3 & $84-116$ & \acs{2SB} & $4-8$ & $8$ & $<39$ & $<43$\\
        4 & $125-163$ & \acs{2SB} & $4-8$ & $8$ & $<51$ & $<82$\\
        5 & $158-211$ & \acs{2SB} & $4-8$ & $8$ & $<55^{9}$ & $<75^{9}$\\
        6v2 & $209-281$ & \acs{2SB} & $4-16^2$ & $24^2$ & $<53^{10}$ & $<66^{10}$\\
        & & & or $4-20^2$ & $32^2$ & & \\
        7v2 & $275-370$ & \acs{2SB} & TBD & $24-32$ & $<72^{10}$ & $<90^{10}$\\
        8v2 & $385-500$ & \acs{2SB} & $4-18^3$ & $28^3$ & $<100^{10,11} $ & $<144^{10}$\\
        & & & & & $<120^{10}$ & \\        
        9 & $602-720$ & \acs{DSB} & $4-12$ & $8(16)^5$ & $<175$ & $<261$\\
        10 & $787-950$ & \acs{DSB} & $4-12$ & $8(16)^5$ & $<230$ & $<344$\\
        \hline
    \end{tabular}
    \begin{minipage}{\textwidth}
    \small
    \vspace*{0.5em}
    \textit{Notes to Table:}     
        \textbf{1.} ``Band 2 Receiver Technical Specifications'' \cite{band2_spe}. 
        \textbf{2.} ``Project Proposal ALMA Band 6v2 Receiver Upgrade: Phase 1'' \cite{band6v2_pp}. Preliminary values pending PDR approval.
        \textbf{3.} ``Project Proposal ALMA Band 8v2 Receiver Upgrade - Phase 1'' \cite{band8v2_pp}. Preliminary values pending PDR approval.
        \textbf{4.} Maximum instantaneous IF bandwidth per polarization.
        \textbf{5.} The maximum bandwidth is approximately doubled in cross-correlation mode with DSB receivers because either sideband can be suppressed using LO offset in each spectral window. 
        \textbf{6.} The maximum specified \ac{SSB} receiver temperatures ($T_{rx}$), averaged over the entire IF band.
        \textbf{7-8.} $T_{rx}$ requirements for the frequency ranges $67-90$~GHz and $90-116$~GHz, respectively.
        \textbf{9.} ``Band 5 Cartridge Technical Specifications'' \cite{band5_spe}.
        \textbf{10.} Table~3 of ``ALMA System Technical Requirements'' (version G) \cite{sys_reqs}.
        \textbf{11.} $T_{rx}$ specified between 390 and 420~GHz \cite{sys_reqs}.
    \end{minipage}
\end{table}
%%%%%%%%%%%%%%%%%%%%%

The \ac{ALMA} receiver bands use different receiver configurations:
%%%%%%%%%%%%%%%%%
\begin{itemize}[itemsep=-0.5em]
    \item Band 1 employs \ac{SSB} receivers, where each polarization provides the full $4-12$~GHz IF bandwidth from the upper sideband only.
    \item Band 2 to 8v2 utilize \ac{2SB} receivers, which provide both upper and lower sidebands separately and simultaneously. Each receiver has four IF outputs—two per polarization. The \ac{IF} bandwidth per output is 4~GHz for Band 3 to 5 and up to 16~GHz for Bands 2, 6v2, 7v2, and 8v2.
    \item Bands 9 and 10 employ \ac{DSB} receivers, where the \ac{IF} contains noise and signals from both sidebands. These receivers have two \ac{IF} outputs, one per polarization, with an \ac{IF} bandwidth of 8~GHz per polarization. \ac{LO} offset can be used to suppress one of the two sidebands. As described in Section~\ref{sec:WSUSignalChain}, this technique shifts correlated signals from the unwanted sideband to a high fringe rate, which effectively spreads them across the spectrum and reduces coherence.
    However, it does not eliminate the noise from the image sideband originating prior to the mixers---this is a fundamental limitation of \ac{DSB} receivers, and neither \ac{LO} offset technique nor 90-degree Walsh phase switching technique (not required for the \ac{WSU}) can overcome this \ac{DSB} noise penalty (it should be noted, however, that for such high frequencies as Bands 9 and 10 that the additional RF hybrid or splitter needed for hardware sideband separation also introduces a non-negligible noise penalty).
\end{itemize}
%%%%%%%%%%%%%%%%%


The \ac{FE} subsystem supports simultaneous operation of up to three bands to prevent overloading the cryostat cooler.
From a hardware point of view, it takes only about $1-2$~seconds to switch between two bands which are powered-on (limited by the \ac{LO} locking and antenna movement).
For bands that are not switched on, it takes around a minute to power up the band, after which a nominal \ac{LO} power must be set. In current \ac{ALMA} operations, approximately 5~minutes is typically allowed for thermal stabilization before a science observation commences, in order to improve operational efficiency. However, the formal technical requirement for thermal stabilization remains 15~minutes for all bands \cite{sys_reqs}.

% RID-5123
During and beyond the \ac{WSU}, the \ac{ALMA} system will operate with a combination of receiver bands upgraded to version~2 (v2) and others remaining as legacy version~1 (v1).
When System Milestone 5 is reached, Bands 6, 7, and 8 will have been upgraded to support wider IF bandwidths and improved sensitivity, while Bands 1, 3, 4, 5, 9, and 10 will continue using legacy cartridges.
ALMA Monitor \& Control software needs to be designed to accommodate heterogeneous Front-End hardware configurations.
Band-specific drivers can be deployed per antenna based on TMCDB definitions (Section~\ref{sec:TMCDB}), ensuring that hardware differences (e.g., tuning algorithms) are correctly managed.
This capability allows simultaneous use of v1 and v2 receiver cartridges within a given band across antennas during the interval over which a band is being upgraded, provided that ALMA control software, databases, and surrounding tools limit the data storage and processing to the v1 \ac{IF} range. 
Finally, to support accurate sensitivity modeling and data reduction, the receiver cartridge version per antenna needs to be recorded in the raw data (\ac{ASDM}) for each receiver band used.


%%%%%%%%%%%%%%%%%%%%%%%%%%%%%%%%%%%%%%%%%%%%%%%%%
\subsection{IF Signal Processing}
\label{sec:WIFP}

%%%%%%%%%%%%%%%%%%%%%%%%%%%%%%%%%%%%%%%%%%%%%%%%%

The \ac{WIFP} is responsible for converting analog \ac{IF} signals from the \ac{FE} into digitized optical signals for transmission via optical fibers.
The \ac{WIFP} consists of three modules:
%%%%%%%%%%%%%%%%%%%
\begin{itemize}[itemsep=-0.5em]
    \item IF Switch for band selection, gain adjustment and slope compensation.
    \item Digitization \& Optical Transmission unit (\ac{DGOT}) for conditioning and converting analog signals to digitized signals and total power detection.
    \item Digital Signal Processing \& Optical Transmission (\ac{DSPOT}) unit for packetizing and transmitting digitized signals over the DTS.
    %\item Time Stamping for precise synchronization of digital samples before transmission to \ac{ATAC} subsystem.
\end{itemize}
%%%%%%%%%%%%%%%%%%%
Accurate time stamping of digitized signals is essential for synchronization across antennas and for achieving phase-coherent correlation.
In the \ac{WSU} architecture, this function is implemented through joint operation of the \ac{DGOT}, and \ac{DSPOT} modules (Section~\ref{sec:TimeStamping}).
The \ac{WIFP} functions are described in detail in \cite{WIFPDescription}.

%%%%%%%%%%%%%%%%%%%%%%%%%%%%%%%%%%%%%%%%%%%%%%%%%
\subsubsection{IF Switch}
%%%%%%%%%%%%%%%%%%%%%%%%%%%%%%%%%%%%%%%%%%%%%%%%%

For the WSU the legacy IF Switch within the FE will be replaced by an equivalent unit meeting the WSU requirements over 2--20 GHz. It is expected to retain the existing IF cables between the receivers and the IF Switch as test data suggest these perform acceptably up to 20 GHz (the cables added to support Band 2 USB, of the same design as the original cables, are all verified up to 20 GHz), however, the bulkhead connectors on the outside of the FE (currently N-type, specified up to 18 GHz) will likely be replaced.


%%%%%%%%%%%%%%%%%%%%%%%%%%%%%%%%%%%%%%%%%%%%%%%%%
\subsubsection{Digitization \& Optical Transmission (DGOT) Unit}
%%%%%%%%%%%%%%%%%%%%%%%%%%%%%%%%%%%%%%%%%%%%%%%%%

The \ac{DGOT} unit, mounted on the cryostat, houses four digitisers, a clock generator and an IF conditioner.
All enclosures are shielded to minimize electromagnetic interference. Signal path inputs are the four analogue IF signals from the \ac{FE}, and the primary outputs are 96$\times$10 Gbps digital optical signals to the \ac{DSPOT}.

The \ac{DGOT} module's IF conditioner provides amplification, slope compensation, fine (0.5dB step) gain adjustment, and filtering of 2--20 GHz (the 2GHz higpass to filter any low frequency spurious signals, and the 20GHz lowpass filter for anti-aliasing). It then splits each \ac{IF} signal into two paths: one for digitization and the other for total power measurements. Detectors measure both full-IF band and narrow-band total power, the latter being filtered through narrow band-pass filters prior to detection.
The full-band detectors are primarily for engineering purposes such as power level optimisation. This is because the full band will have uncertain effective frequency due to gain variation over the bandpass, especially for bands that don't fully cover 2--20 GHz. They will also be affected by inclusion of spurious signals that fall anywhere within the IF (e.g. downconverted YTO harmonics for legacy bands with YTO below 20 GHz). The narrow band detectors support continuum single dish observations, both for antenna calibrations and system characterisation (pointing, focus, beam maps etc.) and Solar mapping. The multiple detector bandpasses sample different frequencies in the IF for wideband receivers, allowing measurement of astronomical flux density distributions, while providing at least one usable bandpass entirely within the IF of each legacy band. The narrow band filters are approximately 2 GHz in bandwidth, similar to the legacy system baseband widths.

Each conditioned IF is fed to a high-speed analog-to-digital converter (ADC) with 6-bit resolution (64 levels) and a sampling rate of 40~GSps. The ADCs receive a  20~GHz sampling clock derived from the 125~MHz reference clock, ensuring synchronization for all digitizers.
The baseline ADC is from Keysight and implements a two-core interleaved architecture with a $\times 2$ demultiplexing stage per core, resulting in 24 output channels per ADC. Consequently, each \ac{DGOT} unit provides 96 output channels, which are transmitted via 96 optical fibers to \acs{DSPOT} module where the signals are captured by FPGA transceivers.



%%%%%%%%%%%%%%%%%%%%%%%%%%%%%%%%%%%%%%%%%%%%%%%%%
\subsubsection{Digital Signal Processing and Optical Transmission (DSPOT) Unit}
%%%%%%%%%%%%%%%%%%%%%%%%%%%%%%%%%%%%%%%%%%%%%%%%%

The \ac{DSPOT} module, enclosed in a shielded housing, is co-located with four 400G-ZR transceivers and an \ac{I/F} board equipped with an \ac{EDFA}, a multiplexer, and a de-multiplexer—all part of the \ac{DTS} subsystem (Section~\ref{sec:DTS}).
The \ac{DSPOT} functions as the Digital Processor Board, responsible for analog-to-digital converter data acquisition at 10~Gbps (96 channels), packetizing digitized \ac{IF} signals, and transmitting them via 400G Ethernet links.
Additionally, the 180-degree Walsh demodulation is implemented in the \ac{DSPOT}.
Each analog-to-digital converter module is assigned a dedicated 400G Ethernet link for reliable high-speed data transmission.

%%%%%%%%%%%%%%%%%%%%%%%%%%%%%%%%%%%%%%%%%%%%%%%%%
\subsubsection{Potentially Combined DGOT+DSPOT Option}
%%%%%%%%%%%%%%%%%%%%%%%%%%%%%%%%%%%%%%%%%%%%%%%%%

At the WIFP PDR the above presented partitioning between DGOT and DSPOT was the preferred option as it reduces some design risk by separating the heat generating FPGAs from the potentially temperature sensitive ADCs. An alternative partitioning with DGOT and DSPOT modules in a single enclosure was also presented. After detailed costing of the WIFP project, the optical interconnections between DGOT and DSPOT have been identified as a major financial cost of the project. The combined DGOT+DSPOT configuration, with direct electrical board-to-board connection between ADCs and FPGAs, thus remains as a potential option due to significant cost saving. There is no difference in functionality between the two options. In case of major issues with the baseline ADC selection, an identified fallback solution would be an integrated ADC-FPGA (Altera Agilex 9 AGRW027), which would only be compatible with combined DGOT and DSPOT functions.


%%%%%%%%%%%%%%%%%%%%%%%%%%%%%%%%%%%%%%%%%%%%%%%%%
\subsubsection{Time Stamping}
%%%%%%%%%%%%%%%%%%%%%%%%%%%%%%%%%%%%%%%%%%%%%%%%%
\label{sec:TimeStamping}

The current \ac{BE} specifications document \cite{BEreqs} defines the \ac{TE} signal requirements, including jitter specifications at each antenna.
The \ac{TE} jitter varies depending on the transmission medium: $\le 0.8$~ns for \ac{LVDS} and $\le 8$~ns for RS422.
However, this level of stability is insufficient for time stamping in the \ac{WSU} architecture, which uses an asynchronous \ac{WSU} \ac{DTS}.

For accurate time stamping of digitized data at an effective sampling rate of 40 GSps, jitter must be $\le 25$~ps \cite{scwg}.
This corresponds to one sampling interval at 40~GSps and ensures that the timing uncertainty remains within the duration of a single sample.
Such precision is essential to maintain proper time alignment and phase coherence across all antennas during correlation.

To achieve this, the \ac{DGOT} regenerates the \ac{TE} signal from \ac{LORR} (delivered in \ac{LVDS} format) using a \ac{DFF} circuit clocked by the 125~MHz reference clock.
This process reduces jitter to $\le 25$~ps, ensuring precise time synchronization within the \ac{WSU} signal chain.
The cleaned new \ac{TE} (nTE) signal is then passed to an \ac{RF} switch, which selects between the standard \ac{IF} signals from the \ac{FE} cartridge and the nTE signal.

When synchronization is required, the \ac{RF} switch routes the nTE signal into the \ac{IF} path within the \ac{DGOT}.
The nTE signal temporarily replaces the \ac{IF} signal, allowing it to be processed identically to an \ac{IF} signal throughout the system.
In \ac{DSPOT}, the \ac{TE} signal is detected within the digitized \ac{IF} sample stream.
A \ac{TE} Detector identifies the \ac{TE} pulse edges and marks the sample where the event occurs.
The \ac{TE} Aligner then shifts the sample stream to ensure that the first sample in each metaframe corresponds to a \ac{TE} event. 
Finally, a \ac{TE}\_index is assigned and inserted into the packet metadata, ensuring accurate synchronization and correlation of data across the system.


%%%%%%%%%%%%%%%%%%%%%%%%%%%%%%%%%%%%%%%%%%%%%%%%%
\subsection{Data Transmission System and Fiber Optic System}
\label{sec:DTS}

%%%%%%%%%%%%%%%%%%%%%%%%%%%%%%%%%%%%%%%%%%%%%%%%%

The \ac{WSU} \ac{DTS} and \ac{FOS} are responsible for transmitting digitized \ac{IF} signals from \acp{WIFP} on at least 66 antennas located at \ac{AOS} to \ac{ATAC} at \ac{OSF}.
%, with infrastructure designed to potentially support up to 70~antennas.
The maximum fiber length between the \ac{AOS} \ac{TB} and an antenna pad is 12.7~km in the current configuration.\footnote{https://confluence.alma.cl/display/ESG/AOS+antenna+pads+optical+fiber+status}
Since \ac{AOS}-\ac{TB} to \ac{OSF} is~30 km, the total fiber distances will be $<43$~km for existing antenna pads and up to 80~km for future antenna pads \cite{dts_pdr_dsn}.

Each antenna transmits four data streams, corresponding to the dual sidebands of dual polarizations per receiver band.
The bit depth and bit rate of each stream are 6-bit and 40~GSps, respectively.
The total raw data rate per antenna (excluding overhead and \ac{FEC}) is 6-bit $\times$ 40~GSps $\times$ 4~streams/antenna $=$ 960~Gbps/antenna.
These functional requirements for the \ac{WSU} \ac{DTS} are summarized in Table \ref{tab:dts_req} \cite{dts_pdr_dsn}.

%%%%%%%%%%%%%%%
\begin{table}[h!]
    \centering
    \caption{\acs{WSU} \acs{DTS} Functional Requirements}
    \begin{tabular}{cc}
         \hline
         Requirement & Specification \\
         \hline
         Number of antennas & at least 66$^1$\\
         Number of optical fiber cables & 2 per antenna between antennas and \ac{ATAC} \\
         Distance& 80~km \\
         Data rate & 960~Gbps/antenna$^2$\\
         Bit-Error Rate & $1\times10^{-15}$ after \ac{FEC}\\
         Data transmission method& Standard Ethernet in full-duplex mode$^3$\\
         \hline
    \end{tabular}
    \begin{minipage}{\textwidth}
    \small
    \vspace*{0.5em}

     \textit{Notes to Table}: 
        \textbf{1.} Processing more than 66 would require additional infrastructure investment, with potential cost and power implications.
        \textbf{2.} The data rate includes only raw data and does not include additional information for data transmission.
        The estimated transmission overhead due to Ethernet packet metaframe, timestamp data, and header formatting is approximately 12\% \cite{atac_icd_dspot}.
        \textbf{3.} Full-duplex communication is standard for Ethernet faster than 1GbE, and neither simplex nor half-duplex communication is supported in relevant IEEE 802.3 standards.
    \end{minipage}
    \label{tab:dts_req}
\end{table}
%%%%%%%%%%%%%%%



%%%%%%%%%%%%%%%%%%%%%%%%%%%%%%%%%%%%%%%%%%%%%%%%%
\subsubsection{WSU DTS Architecture}
%%%%%%%%%%%%%%%%%%%%%%%%%%%%%%%%%%%%%%%%%%%%%%%%%
% RID-4969
The \ac{WSU} \ac{DTS} utilizes 400\,GbE technology based on the \ac{OIF} 400G-ZR standard, which supports long-range transmission of over 80~km.\footnote{The \ac{OIF} 400G-ZR standard specifies a minimum reach of 80~km with \ac{FEC} and allows up to 120~km under favorable conditions, though the \ac{WSU} design baseline assumes a maximum link distance of 80~km.}
Each antenna produces four digitized \ac{IF} signal data streams -- one for each combination of sideband and polarization.
To support this, the \ac{WSU} \ac{DTS} employs four independent 400\,GbE networks per antenna, with one network allocated to each stream.
The use of four independent 400\,GbE networks for the four streams per antenna is an elegant and simple engineering solution, which could lead to stable operation and easy troubleshooting.

At each antenna, four 400G-ZR transceivers (TRXs) transmit these streams as four optical signals.
These signals are multiplexed using \ac{DWDM} technology via a multiplexer (MUX) and then amplified by an \ac{EDFA}.
The combined signal is transmitted to the \ac{OSF} over fiber distances of up to 43~km for existing antenna pads.
The \ac{WSU} \ac{DTS} is designed to support transmission over distances of up to 80~km, which would support future antenna pads aligned with the \ac{ALMA} Development Roadmap mid-term opportunity to implement longer baselines between antennas.

At the \ac{OSF}, the optical signal is demultiplexed (DMUX) into the original four streams and received by another set of 400G-ZR TRXs connected to the \ac{ATAC}.
Digitized signals from the \ac{TP} Array follow the same transmission architecture and are routed through 400\,GbE switches integrated in to the \ac{AVCC} units of \ac{ATAC} (see Section~\ref{sec:TPGS}).

This system architecture is illustrated in Fig.~3.2 of \cite{dts_pdr_dsn}.


%%%%%%%%%%%%%%%%%%%%%%%%%%%%%%%%%%%%%%%%%%%%%%%%%
\subsubsection{Patch Panels at AOS-TB and OSF}
\label{sec:PatchPanel}

%%%%%%%%%%%%%%%%%%%%%%%%%%%%%%%%%%%%%%%%%%%%%%%%%

To support the upgraded \ac{WSU} optical transmission system, a new fiber patch panel will be installed at the AOS Technical Building (AOS-TB) \cite{foc_pp_eu, foc_pp_na}.
This panel will serve as the main interface for connecting new optical fiber cables running between \ac{AOS}-TB and the \ac{OSF}.
 
Additional fiber cables will also be installed to connect the new patch panel with the existing patch panel that currently serves antenna pads.
In this configuration, the new patch panel handles the transmission to \ac{OSF}, while existing patch panel continues to handle connections to the antenna pads.
The new fiber links between the two patch panels bridge the legacy and \ac{WSU} infrastructures, allowing for deployment without observation interruption. 
The \acs{WSU} \ac{FOS} project is responsible for installing the new patch panel, while the \acs{WSU} \acs{DTS} will handle the remaining infrastructure, including floor space and power cabling.

% RID-4971
The new patch panel will be located in the existing patch panel room at AOS-TB, adjacent to the panels currently serving the antenna pads and \ac{ACA} system.
In anticipation of possible future antennas located beyond 15~km from \ac{AOS}-TB, approximately 10 \acp{EDFA} will be installed in the new patch panel racks \cite{dts_pdr_dsn}.

At the \ac{OSF}, a dedicated patch panel for the \ac{WSU} system is installed inside \ac{OCRO} to receive  and distribute the incoming optical signals from \ac{AOS} \cite{foc_pp_na} (Section~\ref{sec:OCRO}).
This panel serves as the primary distribution point at the \ac{OSF} directing the signals to the \ac{DTS} racks installed within \ac{OCRO}. The existing patch panel is a Huber\,+\,Suhner LiSA patch panel system, which is designed for long-term reliability in high-altitude mission-critical environments. Through careful maintenance and a retrofitting performed in 2016, it is expected that their effective lifetime is extended throughout the WSU Program and well into WSU operations.

%%%%%%%%%%%%%%%%%%%%%%%%%%%%%%%%%%%%%%%%%%%%%%%%%
\subsubsection{WSU Fiber Optic System (FOS) between AOS and OSF}
%%%%%%%%%%%%%%%%%%%%%%%%%%%%%%%%%%%%%%%%%%%%%%%%%
\label{sec:FOS}

The \ac{WSU} \ac{FOS} connecting between \ac{AOS} and \ac{OSF} serves as the high-capacity communication infrastructure for the \ac{WSU} system, enabling transmission of digitized \acs{IF} signals over distances of up to 43~km (including the 30~km from the \ac{AOS}-TB to \ac{OCRO}) at the completion of the WSU Program, though the \ac{WSU} system is designed to support distances of up to 80~km \cite{foc_pp_eu} \cite{foc_pp_na}.

The associated cable is a 192-core single-mode glass fiber (9/125~$\mu$m) with an armored jacket, engineered design for direct burial in the environmental conditions of the \ac{ALMA} site. The endpoints will be terminated with meter(s)-long pigtails fitted with high reliability single mode fiber connectors. 
These terminations connect directly to the patch panels at \ac{AOS}-TB and at the \ac{OCRO} within \ac{OSF}, as described in Section~\ref{sec:PatchPanel}.




%%%%%%%%%%%%%%%%%%%%%%%%%%%%%%%%%%%%%%%%%%%%%%%%%
\subsection{Advanced Technology ALMA Correlator (and Beamformer)}
\label{sec:ATAC}

%%%%%%%%%%%%%%%%%%%%%%%%%%%%%%%%%%%%%%%%%%%%%%%%%


The \ac{ATAC} subsystem receives digitized data with accurate time-stamping from each \ac{ALMA} antenna’s \ac{WIFP} subsystem via \ac{DTS}, along with delay models and observing configurations from Control. \ac{ATAC} will produce Imaging Correlation (IC) visibilities for the 7-m and 12-m Arrays, with a native channel bandwidth of 13.5~kHz for Standard Interferometry observing modes while optionally producing Zoom Windows (with even higher spectral resolution traded for bandwidth) or \ac{VLBI} (phased-array) 
Beams.
\ac{ATAC} is designed to produce up to 32~GHz per polarization of correlated data and \ac{VLBI} beam channels, with 16~GHz per polarization output available in its initial deployment, and a straightforward expansion path to 32~GHz per polarization with additional funding. 
Table~\ref{tab:atac_req} summarizes the key requirements of the \ac{ATAC} subsystem. Table~\ref{tab:atac_CH_Tint} provides the finest channelization and integration time options for visibilities, which are configurable per spectral window.

%%%%%%%%%%%%%%%%%%
\begin{table}[h]
    \centering
    \caption{\ac{ATAC} Key Requirements \cite{atac_sdd_dd-1} \cite{atac_spe_ds-1}}
    \begin{tabular}{lc}
         \hline
         Requirement & Specification \\
         \hline
         Maximum \ac{CBW} & 16 (32)$^1$ GHz per polarization \\
         Maximum VLBI/Beamforming Bandwidth & 16 (32)$^1$ GHz per polarization$^2$\\
         Maximum Number of Spectral Windows & 80 (160)$^1$ \\
         Number of Channels (per pol) & 80 $\times$ 14880 (160 $\times$ 14880)$^1$ \\
         Native Channel Width & 13.5~kHz \\
         Bandwidth with Zoom & 16~GHz/ZF (32~GHz/ZF)$^{1,2,3}$ \\
         Spectral Channel Width with Zoom & 13.5~kHz$^2$\\
         \hline
    \end{tabular}
    \begin{minipage}{\textwidth}
    \small
    \vspace*{0.5em}
        \textit{Notes to Table:} 
        \textbf{1.} Initial WSU deployment 2$\times$ BW  (after future expansion project to 4$\times$ BW). 
        \textbf{2.} Assuming all processing resources are in \ac{ICV} Mode, i.e. no Zoom Windows are employed. 
        \textbf{3.} Where the Zoom Factor (ZF) is 2, 4, or 8, and all processing resources are in \ac{ICZoom} Mode.
    \end{minipage}
    \label{tab:atac_req}
\end{table}
%%%%%%%%%%%%%%%%%%

%%%%%%%%%%%%%%%%%%
\begin{table}[h]
    \centering
    \caption{ATAC Minimum Visibility Channelization and Integration Time \cite{atac_sdd_dd-1}}
    \begin{tabular}{cccc}
         \hline
         Correlation Mode &  Minimum &  Minimum & Science Bandwidth \\
          &  Channel Bandwidth & Integration Time & per Polarization (GHz)$^1$ \\
         \hline
         Native Spectral Resolution &  13.5~kHz& 1024~ms & 16 (32) \\
         Fast Integration &  1.08~MHz & 16~ms & 16 (32) \\
         Zoom Spectral Resolution & 1.6875~kHz & 1024~ms & 1 (2) \\
         \hline
    \end{tabular}
    \begin{minipage}{0.96\textwidth}
    \small
    \vspace*{0.5em}
        \textit{Notes to Table:} 
        \textbf{1.} Initial WSU deployment 2$\times$ BW (after future expansion project  to 4$\times$ BW).
    \end{minipage}
    \label{tab:atac_CH_Tint}
\end{table}
%%%%%%%%%%%%%%%%%%


\noindent Additional highlights of the \ac{ATAC} design include:

%%%%%%%%%%%%%%%%%%%%%
\begin{itemize}[itemsep=-0.5em]
    \item Ability to ingest 40.0~GSps per sideband per polarization per antenna from at least 66 antennas,\footnote{While the \ac{ATAC} subsystem is required to process digitized data from at least 66~antennas the design may accommodate a modest increase in antenna count -- e.g., up to 72 (TBC) -- if needed. This aligns with the ALMA Development Roadmap mid-term opportunity to increase the number of 12-m antennas.} with support for up to eight independent concurrent correlator subarrays.
    \item An FFX architecture based on polyphase filter banks with a coarse channelization stage that generates \ac{FS} and a second stage 16~k channelizer for fine channelization prior to correlation that yields all four polarization products.
    \item All-digital delay and phase tracking with high accuracy for up to 35~km baseline lengths.
    \item A digital processing efficiency of at least 0.995. 
    \item In addition to Imaging Correlation visibilities, \ac{ATAC} can also optionally produce:
   %%%%%%%%
   \begin{itemize}[itemsep=-0.5em]
%     \item On a per \ac{FS} basis, Zoom Spectral Windows that trade bandwidth for higher spectral resolution by factors of 2, 4, 8, or 16; or
     \item On a per \ac{FS} basis, Zoom Spectral Windows that trade bandwidth for higher spectral resolution by factors of 2, 4, or 8; or
     \item One \ac{VLBI} beam per correlator subarray at full science bandwidth (with a maximum of two such subarrays at any given time per \ac{FS}) using true-delay beamforming. A \ac{VLBI} beam can be formed from 1 to 66 antennas. For high spectral resolution applications, an optional 128~MHz or 16~MHz beam-channel can also be produced. 
   \end{itemize}
   %%%%%%%
\noindent These options are accessed using one of the two \ac{ATAC} Modes: \ac{ICZoom} or \ac{ICV}, which can be independently configured for each \ac{FS}.

\item Provides \ac{ATAC} functionality within the \ac{WSU} \ac{HILSE}, allowing tests using either live digitized sky signals from up to four selected ALMA antennas or simulated digital signals (Section~\ref{sec:ATACHILSE}).
\acs{ATAC}.\acs{HILSE} operates independently from the operational correlator and supports both \ac{ATAC} modes (\ac{ICZoom} and \ac{ICV}).

\end{itemize}
%%%%%%%%%%%%%%%


The \ac{ATAC} subsystem is located in \ac{OCRO} (Section~\ref{sec:OCRO}), and consists of the ATAC \ac{CBF}, ATAC Monitor \& Control, \ac{CDP}, Configuration Library (CONFIG\_LIB), ATAC \ac{HILSE}, and Rack, Power, and Cooling (RPC).
Their functionalities are described below, see \cite{atac_sdd_dd-1} for more detail. 


%%%%%%%%%%%%%%%%%%%%%%%%%%%%%%%%%%%%%%%%%%%%%%%%%
\subsubsection{ATAC Processing Stages: ATAC.CBF (AVCCs \& AFSPs) and ATAC.CDP}
%%%%%%%%%%%%%%%%%%%%%%%%%%%%%%%%%%%%%%%%%%%%%%%%%

The \ac{ATAC} subsystem can produce a continuous set of spectral channels covering up to 16~GHz per polarization (with the potential to be expanded to up to 32~GHz per polarization in the future), with no gaps, aliasing, or correlator-induced changes in sensitivity.
\ac{ATAC} processing is done through three stages as shown in Fig.~\ref{fig:ATACFig1} and described in detail in \cite{atac_sdd_dd-1}.

%%%%%%%%%%%%%%%%%%%%
\begin{figure}[htb]
\includegraphics[width=\textwidth]{Figures/ATACFig1.png}
\caption{Block diagram showing the AVCC/AFSP/CDP split. \label{fig:ATACFig1}} 
\end{figure}
%%%%%%%%%%%%%%%%%%%%

\noindent
\textbf{Stage~1: \ac{ALMA} Very Coarse Channelizers (\acs{AVCC}s)} perform the first "F" stage common to \ac{ICZoom} and \ac{ICV} \ac{ATAC} modes.
This stage utilizes an \acs{OSPFB} to produce oversampled \acp{FS} that are 222.222~MS/s wide per polarization, with the central 200.88~MHz (14880~channels) containing un-aliased, science-quality data.
Each \ac{AVCC} generates 162~dual-polarization \acp{FS} (81 per sideband) to cover 32~GHz of science bandwidth.
These 81~\acp{FS} per sideband can be selected from any 16~GHz segment within the 20~GHz of digitized bandwidth to accommodate receivers with different IF bandwidth properties (Fig.~\ref{fig:ATACFig2}).
Each \ac{FS} can be transmitted to one or more \acp{AFSP} for further processing.


\begin{figure}[htb]
\centering
\includegraphics[width=0.92\textwidth]{Figures/ATACFig2.png}
\caption{Demonstration of channelization stages. The top row shows how two potential \ac{2SB} receiver IF configurations would map to the $0-20$~GHz of IF bandwidth that will be digitized for each sideband at the antennas. The \ac{AVCC} stage only shows one of the two sidebands that are channelized in this stage. The \ac{AFSP} stage produces nearly independent channels \cite{atac_sdd_dd-1}. \label{fig:ATACFig2}} 
\end{figure}
%%%%%%%%%%%%%%%%%%%%

\noindent
\textbf{Stage 2: ALMA Frequency Slice Processors (\acs{AFSP}s)} provide delay and phase tracking, removal of LO Offset per antenna, fine channelization, corner turn, correlation and create full polarization visibility sets. Each \ac{AFSP} processes one \ac{FS} for all 66~\ac{ALMA} antennas\footnote{When the antennas are split into multiple Arrays, the same \ac{AFSP} may be processing a different \ac{RF} range on the different sets of antennas corresponding to the different correlator subarrays.} and can be configured to run in one of two modes: \ac{ICZoom} or \ac{ICV}, which in addition to the Imaging Correlation functions can optionally produce Zoom Windows or \ac{VLBI} Beamforming, respectively.
\ac{VLBI} beams are sent directly to the \ac{VLBI} recorders,\footnote{The current VLBI recorders are not part of the ATAC subsystem and cannot record more than 8~GHz of bandwidth per polarization, summed over all correlator subarrays. The full bandwidth improvements enabled by ATAC will require future upgrades to these recorders.} while correlated visibilities are sent to the \ac{CDP} for further processing.
Each \ac{AFSP} can only run one mode at a time but can be reconfigured to match observation needs.
Additional details regarding how \ac{ICZoom} and \ac{ICV} can be used concurrently, see \cite{Memo023}.

\noindent
\textbf{Stage 3: Correlator Data Processor (\acs{CDP})} receives visibility sets from \acp{AFSP} operating in \ac{ICZoom} or \ac{ICV}.
The \ac{CDP} seamlessly combines multiple, adjacent \acp{FS} together to cover the contiguous
bandwidth of each spectral window requested by the user; in the initial deployment up to 80 independent spectral windows can be created (i.e. one \ac{FS} per spectral window).
For each spectral window, the \ac{CDP} performs the user-requested channel and time averaging of visibility data to achieve the required spectral resolution and integration time.
The \ac{CDP} also applies the Van Vleck corrections on visibilities, normalizes cross-correlation visibilities based on the power spectra of the corresponding autocorrelations, creates the channel average datastream for each spectral window, generates \ac{BDF} files and transmits them to TelCal and Archive.
An innovative aspect of the \ac{ATAC} subsystem design is that the bulk of \ac{CDP} processing will occur on the \acp{CPU} that are co-located with the \ac{AVCC} and \ac{AFSP} \ac{FPGA} hardware in rack-mounted servers called TeraBoxes \cite{atac_sdd_dd-1}.


%%%%%%%%%%%%%%%%%%%%%%%%%%%%%%%%%%%%%%%%%%%%%%%%%
\subsubsection{ATAC Configuration Library}
%%%%%%%%%%%%%%%%%%%%%%%%%%%%%%%%%%%%%%%%%%%%%%%%%

The \ac{ATAC} Configuration Library supplies an Application Programming Interface to interactively build valid Correlator Configurations including the signal processing resources needed, Spectral Window properties, visibility time averaging, channel averaging, polarization products, etc.
The Configuration Library will be callable by any \ac{ALMA} subsystem such as the \ac{OT} and, internally \acs{ATAC}.M\&C and \acs{ATAC}.\acs{CDP}.


%%%%%%%%%%%%%%%%%%%%%%%%%%%%%%%%%%%%%%%%%%%%%%%%%
\subsubsection{ATAC Monitor \& Control}
%%%%%%%%%%%%%%%%%%%%%%%%%%%%%%%%%%%%%%%%%%%%%%%%%
\label{sec:ATAC.M&C}

\ac{ATAC} Monitor \& Control (\ac{ATAC}.M\&C) software is an integral part of the \ac{ATAC} subsystem which can be functionally described as a “gateway” for all monitor and control requests originating outside of it.
The \ac{ATAC}.M\&C provides:
%%%%%%%%%%%%%%%%%%
\begin{itemize}[itemsep=-0.5em]
    \item A centralized point of control for all \ac{ATAC} hardware including signal processing equipment as well as critical support equipment such as power and cooling systems.
    \item Translation of ACS commands from external subsystems to commands used internally within the \ac{ATAC}.M\&C software.
    \item Routing of \ac{ATAC} internal Alarms, events, and logs to the appropriate \ac{ACS} components.
    \item Reporting of summarized status for all equipment and correlator subarrays contained in \ac{ATAC}.
\end{itemize}
%%%%%%%%%%%%%%%%%%
See \cite{atac_sdd_dd-1} and \cite{atac_spe_ds-1} for additional details on \ac{ATAC}.M\&C.

%%%%%%%%%%%%%%%%%%%%%%%%%%%%%%%%%%%%%%%%%%%%%%%%%
\subsubsection{ATAC HILSE}
%%%%%%%%%%%%%%%%%%%%%%%%%%%%%%%%%%%%%%%%%%%%%%%%%
\label{sec:ATACHILSE}

\ac{ATAC} \ac{HILSE} is a standalone system within the \ac{ATAC} subsystem, designed to provide a representative slice of \ac{ATAC} for testing \ac{ALMA} hardware, firmware, and software in a separate \ac{APE} without disrupting science operations (see Section~\ref{sec:online}).
It offers the following key capabilities:
%%%%%%%%%%%%%%%%%
\begin{itemize}[itemsep=-0.5em]
    \item Simulates wideband digitized data from the \ac{WIFP} subsystem, supporting up to four antennas with continuous, real-time data for long integrations.
    \item Ingests real digital signals from up to four ALMA antennas, selectable via the \ac{AOS} Patch Panel.
    \item Allows selection between simulated test vectors and real \ac{ALMA} antenna data for testing.
    \item Includes sufficient portions of \ac{ATAC}.\ac{CBF}, \ac{ATAC}.M\&C, and \ac{ATAC}.\ac{CDP} hardware to run operational software and \ac{FPGA} firmware independently.
    \item Connects to standalone instances of Control, TelCal, Archive, and \ac{VLBI} recorders.
\end{itemize} 
%%%%%%%%%%%%%%%%
The capabilities and architecture of \ac{ATAC} \ac{HILSE} are described in detail in \cite{atac_sdd_dd-1}.

%%%%%%%%%%%%%%%%%%%%%%%%%%%%%%%%%%%%%%%%%%%%%%%%%
\subsubsection{ATAC Rack, Power, Cooling}
%%%%%%%%%%%%%%%%%%%%%%%%%%%%%%%%%%%%%%%%%%%%%%%%%

The \ac{ATAC} subsystem will employ liquid cooling to significantly decrease the amount of air-cooling needed for its processing hardware.
\ac{ATAC} will provide the secondary liquid cooling network (``back of rack'' manifolds and cooling lines that circulate liquid in the TeraBox servers) that connects to the \ac{OCRO} primary liquid cooling network, which includes an outdoor heat exchanger (also see \ref{sec:OCRO}). 


%%%%%%%%%%%%%%%%%%%%%%%%%%%%%%%%%%%%%%%%%%%%%%%%%
\subsection{Total Power GPU Spectrometer (TPGS)}
\label{sec:TPGS}

%%%%%%%%%%%%%%%%%%%%%%%%%%%%%%%%%%%%%%%%%%%%%%%%%

The \ac{TPGS} subsystem employs the same FFX-type signal processing chain as the \ac{ATAC} subsystem (Section~\ref{sec:ATAC}), ensuring full functional compatibility between the two subsystems.
The \ac{TPGS} subsystem is designed to provide 32~GHz \ac{IF} bandwidth per polarization capability and generates auto- and cross-correlation products for four 12-m antennas of the \ac{TP} Array, with the cross-correlation products used solely for calibration purposes such as focus and pointing. Table~\ref{tab:tpgs_req} outlines the key requirements of the \ac{TPGS} subsystem \cite{TPGS_cosdd}.

%%%%%%%%%%%%%%%%%%%%%%%%%
\begin{table}[h]
    \centering
    \caption{\ac{TPGS} Key Requirements \cite{TPGS_cosdd}}
    \label{tab:tpgs_req}
    \begin{tabular}{lc}
         \hline
         Requirement & Specification \\
         \hline
         Maximum \ac{CBW} & 32 GHz per pol \\
         Maximum number of spectral windows & 160 \\
         Number of channels (per pol)$^1$ & 160 $\times$ 14880 \\
         Channel width at Max \ac{CBW}$^1$ & 13.5~kHz \\
         Finest channel width$^1$ & 13.5 kHz / 8 (1.6875~kHz) \\
         \hline
    \end{tabular}
    \begin{minipage}{0.68\textwidth}
    \small
    \vspace*{0.5em}
        \textit{Notes to Table:}
        \textbf{1.} for dual and full polarization products.
    \end{minipage}
\end{table}
%%%%%%%%%%%%%%%%%%%%%%%%%

Fig.~\ref{fig:TPGSFig1} illustrates the schematic diagram of the TPGS subsystem \cite{TPGS_cosdd}. The \ac{TPGS} subsystem consists of two primary software components: \ac{TPGS}-\ac{PU} and \ac{TPGS}-\ac{CU}. The \ac{TPGS}-\ac{PU} hardware comprises \ac{GPU} servers, with each \ac{GPU} server dedicated to processing one of the two sidebands independently. These servers perform tasks such as coarse channelization, resampling, delay correction, \ac{LO} offset removal, fine channelization (including zoom modes), and calculating auto- and cross-correlations for their respective sidebands. The \ac{TPGS}-\ac{CU}, which runs on a \ac{MCF} server, acts as the interface between Control (ALMA Control Software) and the \ac{TPGS}-\ac{PU}.

Additionally, the \ac{TPGS}-\ac{CU} is responsible for quantization correction, stitching multiple processed \acp{FS}, generating \ac{BDF} files for each spectral window, and transmitting these \acp{BDF} to the Archive and TelCal subsystems.

%%%%%%%%%%%%%%%%%%%%%%%%%
\begin{figure}[htb]
\centering
\includegraphics[width=\textwidth]{Figures/TPGS_ConfigurationOption1.png}
\caption{Schematic diagram of the \ac{TPGS} subsystem. The TPGS subsystem consists of a general-purpose server for \ac{MCF}, two \ac{GPU} servers for data processing, and one network switch for internal communication.
Additionally, four \ac{AVCC} units and one \ac{CDP} 400 GbE switch are part of the \ac{ATAC} subsystem. Each \ac{AVCC} unit includes a TeraBox and a 32-port 400 GbE network switch. Each \ac{GPU} server requires either twelve 200 GbE ports or eight 400 GbE ports for connectivity \cite{TPGS_cosdd}.\label{fig:TPGSFig1}} 
\end{figure}
%%%%%%%%%%%%%%%%%%%%%%%%%

The \ac{TPGS} \ac{HILSE}, a test environment for the \ac{TP} Array, will be integrated into the \ac{ALMA} \ac{WSU} \ac{HILSE} via the \ac{ATAC} \ac{HILSE}.
The \ac{TPGS} \ac{HILSE} consists of one \ac{MCF} server and one \ac{GPU} server.
These components perform the same functions as the \ac{TPGS}-\ac{CU} and \ac{TPGS}-\ac{PU} but are designed to handle only 16~GHz \ac{IF} bandwidth per polarization for a single sideband \cite{TPGS_cosdd}.

All \ac{TPGS} subsystem equipment, including the \ac{TPGS} \ac{HILSE}, will be installed in two racks in the \ac{OCRO} (Section~\ref{sec:OCRO}).
The system comprises two \ac{MCF} servers, three \ac{GPU} servers, two network switches, and two \acp{PDU}. 



%%%%%%%%%%%%%%%%%%%%%%%%%%%%%%%%%%%%%%%%%%%%%%%%%
\subsection{Local Oscillators}
\label{sec:LO}

%%%%%%%%%%%%%%%%%%%%%%%%%%%%%%%%%%%%%%%%%%%%%%%%%

The \ac{ALMA} local oscillators play a crucial role in time synchronization across the array, generating sinusoidal signals necessary for converting observed sky frequencies to intermediate frequencies ($2-20$~GHz) for the transmission of digitized signals to \ac{ATAC} and \ac{TPGS}.
Additionally, they provide precise frequency references to achieve the desired sky frequency and interferometer phase, incorporating features such as Walsh phase switching, and other interferometer-specific features \cite{BE_dsn} \cite{Timing_alma}.
The \ac{WSU} signal chain continues to rely on the existing \ac{PLO} subsystem to provide stable and coherent reference signal distribution across the array.
No architectural changes are required in the \ac{PLO} subsystem to support the \ac{WSU} system architecture (Section~\ref{sec:RefGen}).


%%%%%%%%%%%%%%%%%%%%%%%%%%%%%%%%%%%%%%%%%%%%%%%%%
\subsubsection{Reference Generation and Distribution}
%%%%%%%%%%%%%%%%%%%%%%%%%%%%%%%%%%%%%%%%%%%%%%%%%
\label{sec:RefGen}

% RID-4981/4897/4961/4883
The \ac{PLO} subsystem generates and delivers high stability reference signals to the \ac{ALMA} antennas to ensure that all antennas of the array operate coherently and synchronously with each other and effectively make the antenna array collectively operate as a giant single receiving antenna with unprecedented resolution and sensitivity in the mm and sub-mm range.
The \ac{PLO} subsystem functions are mainly implemented by the \ac{CLOA}, \ac{LPR}, and \ac{LORR}.

The \ac{CLOA} generates the \ac{LO1} photonic reference (Section~\ref{sec:LO1}) as well as the low frequency references (48~msec \ac{TE}, 5~MHz, 125~MHz, and 2~GHz) and distributes them to all antennas and other subsystems that require direct reference signals, based on the \ac{MFS} -- excluding \ac{ATAC} and \ac{TPGS} subsystems, which recover the \ac{TE} signal from the digitized data stream via the \ac{WIFP} subsystem.
The \ac{CLOA} also includes the \acp{LLC}, which compensate for thermal expansion and mechanical stresses on the optical fibers between the \ac{CLOA} and the \ac{FE} assembly located at each antenna during each \ac{EB}.
This compensation is achieved by measurement of the round trip phase using the laser output from the \ac{ML}, which operates in parallel with the distribution of the \ac{LO1} photonic reference to the antennas. 
No changes are needed in this aspect of the \ac{LO1} subsystem to support the \ac{WSU}.

To synchronize the 48~msec \ac{TE} signal with absolute \ac{UTC}, a rack-mounted \ac{GPS} receiver located in the computer room of the \ac{AOS}-TB provides a 1~\ac{PPS} signal (see Fig.~1 of \cite{Timing_alma}). 
This signal is distributed to both the \ac{CRD} and the \ac{MFS} (hydrogen maser), which are located in the \ac{CLOA}.
At system initialization, the \ac{CRD} uses the 1\ac{PPS} to align the \ac{TE} signal with \ac{UTC}.
The \ac{CRD} then counts 125~MHz pulses internally to maintain time, enabling the \ac{ARTM} to calculate and maintain \ac{AT} based on the number of elapsed \ac{TE} intervals since initialization.
This mechanism ensures that the \ac{TE} signal, which is distributed through the Array, remains synchronized with \ac{UTC}, with one \ac{TE} pulse occurring every 48~msec and a guaranteed alignment with the 1\ac{PPS} every 6~seconds.
The system also allows for comparison between the 1\ac{PPS} signal from the \ac{GPS} located in the correlator room and the \ac{MFS}, enabling tracking of the long-term drift between the \ac{GPS}-based UTC and the maser-based time.

In parallel, absolute \ac{UTC} time is distributed to real-time computers via an \ac{NTP} server currently located at the \ac{OSF}.
These computers synchronize their clocks at boot time using \ac{NTP} client requests.
Over the next year or so, a new integrated \ac{GPS}/\ac{NTP} server unit (Brandywine M212) will be installed at the \ac{AOS} to replace the current configuration of three \ac{GPS} units at \ac{AOS} plus one \ac{NTP} server at \ac{OSF}.


%%%%%%%%%%%%%%%%%%%%%%%%%%%%%%%%%%%%%%%%%%%%%%%%%
\subsubsection{First Local Oscillator (LO1)}
%%%%%%%%%%%%%%%%%%%%%%%%%%%%%%%%%%%%%%%%%%%%%%%%%
\label{sec:LO1}

In the \ac{CLOA}, the \ac{ML} operates at a fixed wavelength of 1556~nm, while the tunable \ac{SL} is offset-locked to the \ac{ML}.
The offset frequency can be set anywhere within the ranges of $27-38$~GHz and $65-122$~GHz.
The outputs of the \ac{ML} and \ac{SL} are combined to produce an optical beat signal, which is transmitted to the antennas via optical fiber.
At the antenna, the optical beat signal is amplified by \ac{LPR}, located in the FE. 
The amplified signal is then directed to the photomixer in the \ac{WCA} of the active receiver band. 
The photomixer generates the desired \ac{LO1} reference signal as the difference frequency between the two phase-locked optical tones.

The photonic \ac{LO1} reference signal is distributed to all antennas within a given array. 
At each antenna, a \ac{FLOOG} applies an individual frequency and phase offset to the common reference before feeding it into the \ac{PLL} circuit that drives the \ac{YTO}. 
The \acp{FLOOG} of all antennas are continuously tracked throughout an observation and serves multiple functions:
%%%%%%%%%%%%%%%%%%%%%%
\begin{itemize}[itemsep=-0.5em]
    \item \ac{LO} offset: Applies a small frequency offset to the \ac{LO1} signal to mitigate internally generated interference and suppress correlated signals from the image sideband in \ac{DSB} receivers. This offset is removed during digital processing at \ac{ATAC} and \ac{TPGS} \cite{atac_sdd} \cite{TPGS_cosdd}. The signal-processing implications of this technique are described in Section~\ref{sec:WSUSignalChain}.
    \item Walsh phase switching: Modulates the \ac{LO1} phase using 180-degree Walsh functions to suppress spurious signals and DC offsets. Demodulation is performed in the WIFP’s DSPOT module \cite{WIFPDescription}. The 90-degree phase switch previously used for sideband separation is no longer supported in the \ac{WSU} architecture.
\end{itemize}
%%%%%%%%%%%%%%%%%%%%%%%%%%%

In the \ac{WCA}, a \ac{YTO} serves as the fundamental oscillator, generating a periodic carrier signal. 
The \ac{YTO} frequency is band-dependent, ranging from approximately 13.76 to 40.18~GHz for for Bands 1 to 10 (Table~\ref{tab:lo1_spec}).
This carrier signal is multiplied up for Bands 2 to 10, and phase locked to the corresponding \ac{LO1} photonic reference. 
This signal is then subsequently multiplied up to the desired \ac{LO1} frequency for Bands 2 and 4 through 10.
%This carrier signal is phase-locked to the \ac{LO1} photonic reference and then multiplied up to the desired \ac{LO1} frequency in Bands 2 to 10. 
However, in Band 1, the \ac{YTO} signal is directly phase-locked to the \ac{LO1} photonic reference without requiring multiplication. Table~\ref{tab:lo1_spec} summarizes the \ac{LO1} frequency generation across all bands.

%%%%%%%%%%%%%%%%%%%%%%%%%
% RID-4936
%%%%%%%%%%%%%%%%%%%%%%%%%
\begin{table}[hbt]
    \centering
    \caption{First Local Oscillator Specifications}
    \label{tab:lo1_spec}
    \begin{tabular}{llllll}
        \hline
         Band & LO range (GHz) & YTO range & $N_{WCA}^6$ & \ac{WCA} LO output & $N_{multiplier}^7$\\
         & & (GHz) & & frequency (GHz) & \\
         \hline
         1 & $31-38^2$ & $30.82-40.18$$^4$ & $-$ & $30.8-40.2$ & $-$\\
         2 & $79-104$ & $19.63-26.12$$^4$ & $4$ & $79.0-104.0$ & $-$\\
         3 & $92-108$ & $15.27-18.05$$^4$ & $6$ & $92.0-108.0$ & $-$\\
         4 & $133-155$ & $22.08-25.90$$^4$ & $3$ & $66.5-77.5$ & $2$\\
         5 & $166-203^3$ & $13.76-16.98$$^4$ & $6$ & $83.0-101.5$ & $2$\\
         6v2 & $221-265$ & $24.56-29.44^5$ & $3$ & $73.7-88.3^5$ & $3$\\
         7v2 & $283-365$ & TBD & TBD & TBD & TBD \\
         8v2 & $393-492$ & $21.78-27.33$ & $3$ & $65.3-82.0$ & $6$\\
         9$^1$ & $610-712$ & $22.51-26.45$$^4$ & $3$ & $67.78-79.11$ & $9$\\
         10$^1$ & $795-942$ & $14.66-17.49$$^4$ & $6$ & $88.33-104.67$ & $9$\\
         \hline
    \end{tabular}
    \begin{minipage}{\textwidth}
    \small
    \vspace*{0.5em}
    \textit{Notes to Table:}
        \textbf{1.} The LOs were designed to cover the range required by a possible $4-8$~GHz 2SB receiver.
        \textbf{2.} The LO range upper limit is 40~GHz, but the laser synthesizers only reach 38~GHz.
        \textbf{3.} ``Band 5 Cartridge Technical Specifications'' \cite{band5_spe}
        \textbf{4.} ``ALMA System Block Diagram`` \cite{alma_sbd}
        \textbf{5.} ``Project Proposal ALMA Band 6v2 Receiver Upgrade: Phase 1'' \cite{band6v2_pp}.
        \textbf{6.} Harmonic multiple in a \ac{WCA}.
        \textbf{7.} Harmonic multiple by multipliers in a \ac{CCA}.
    \end{minipage}
\end{table}
%%%%%%%%%%%%%%%%%%%%%%%%%

The broader \ac{IF} range of the \ac{WSU} ($2-20$~GHz) now overlaps with certain YTO frequencies, as shown in Table~\ref{tab:lo1_spec}, although the original design placed \ac{YTO} frequencies above the \ac{ALMA} legacy \ac{IF} range ($4-12$~GHz).
To address this issue, mitigation strategies should be considered for future \ac{WCA} upgrades.
This overlap may result in spurious mixing products or increased \ac{RFI} if \ac{YTO} harmonics fall within the \ac{WSU} \ac{IF} range or the \ac{RF} ranges of receivers.
Measures such as adjusting the \ac{YTO} parking frequency based on the location of active spectral windows should be considered to minimize \ac{RFI} risks.


%%%%%%%%%%%%%%%%%%%%%%%%%%%%%%%%%%%%%%%%%%%%%%%%%
\subsection{Site Infrastructure}
%%%%%%%%%%%%%%%%%%%%%%%%%%%%%%%%%%%%%%%%%%%%%%%%%

This section outlines the key infrastructure upgrades required to support the \ac{WSU} at the \ac{OSF} and \ac{SCO}. 
It focuses on the development of the new \ac{OCRO} and the implications for electrical power provisioning during the parallel deployment phase.
These upgrades are essential to ensure that the \ac{ATAC}, \ac{TPGS}, \ac{DTS}, and associated subsystems can be installed and operated reliably without affecting ongoing observatory operations


%%%%%%%%%%%%%%%%%%%%%%%%%%%%%%%%%%%%%%%%%%%%%%%%%
\subsubsection{OSF Correlator Room (OCRO)}
\label{sec:OCRO}
%%%%%%%%%%%%%%%%%%%%%%%%%%%%%%%%%%%%%%%%%%%%%%%%%

The current control room in the \ac{OSF} Technical Building will be refurbished and transformed into the \ac{OCRO}, a dedicated infrastructure supporting the operational requirements of \ac{ATAC} (Section~\ref{sec:ATAC}), \ac{TPGS} (Section~\ref{sec:TPGS}), the \ac{DTS} fiber electronics (Section~\ref{sec:DTS}), and potentially future computing systems  \cite{OCRO_DSN}.
The refurbishment includes the installation of new power distribution lines with \ac{UPS} support, the deployment of both air-based and liquid-based cooling systems, seismic rack reinforcements, controlled access, and a fire protection system.

The \ac{OCRO} includes an Entrance Area, Service Area, Room Area, Electrical Trench, Cooling Trench, Chiller Area, and UPS Room (Fig.~\ref{fig:OCRO-layout}).
The Room Area will accommodate 28~racks in two rows: 22~racks for \ac{ATAC} (including \ac{HILSE}), 1~rack for \ac{TPGS}, 2~racks for \ac{DTS}, and 3~spare racks allocated for potential future systems.
The facility will include \ac{CRAH} units and \acp{CDU} for cooling, and the cooling system will utilize both air and liquid cooling methods. Electrical and liquid cooling distribution will be routed through dedicated underfloor trenches.
A fire protection system using FK-5-1-12 clean agent fire suppressant, along with a \ac{VESDA} smoke detection system, will be installed with redundancy.


%%%%%%%%%%%%%%%%
\begin{figure}[hbt]
    \centering
    \includegraphics[width=\textwidth]{Figures/Figure8_OCRO.png}
    \caption{Layout of the new OSF Correlator Room \cite{OCRO_DSN}.}
    \label{fig:OCRO-layout}
\end{figure}
%%%%%%%%%%%%%%%%


%%%%%%%%%%%%%%%%%%%%%%%%%%%%%%%%%%%%%%%%%%%%%%%%%
\subsubsection{OSF and SCO Control Rooms}
\label{sec:controlRooms}
%%%%%%%%%%%%%%%%%%%%%%%%%%%%%%%%%%%%%%%%%%%%%%%%%
Control rooms are established in both the \ac{SCO} and \ac{OSF}, providing redundant and functionally equivalent infrastructure to operate \ac{ALMA} efficiently.
Each control room is equipped with two sets of consoles to support concurrent subarray operation modes. 
Additional monitors display the status of ancillary systems essential for scientific observations, such as weather stations, weather forecasts for the Chajnantor Plateau, surveillance cameras, etc.

Both control rooms are capable of supporting full array operations. Routine daytime operations are primarily conducted from the \ac{SCO} control room, while nighttime observations are currently carried out mainly from the \ac{OSF} control room.
The \ac{OSF} control room also ensures operational continuity in the event of major failures, such as disruption in the fiber link between \ac{OSF} and \ac{SCO}, and serves as the main interface point for array subsystems during assembly, integration, verification, and commissioning phases -- both for the current \ac{ALMA} and \ac{WSU} systems.

To accommodate the construction of the \ac{OCRO}, the original OSF control room has been dismantled.
A new \ac{OSF} control room has been constructed at a nearby location as part of the \ac{CRR} project. 
This new facility is designed to match the size and functionality of the \ac{SCO} control room and continues to support all operational roles of the former \ac{OSF} control room, including backup control, nighttime operations, and subsystem interfacing. 
The layout of the \ac{SCO} control room is shown in Fig.~\ref{fig:SCO-Control-Room-layout}.

%%%%%%%%%%%%%%%%
\begin{figure}[ht!]
    \centering
    \includegraphics[width=\textwidth]{Figures/SCO-Control-Room.png}
    \caption{Layout of the SCO Control Room.}
    \label{fig:SCO-Control-Room-layout}
\end{figure}
%%%%%%%%%%%%%%%%


%%%%%%%%%%%%%%%%%%%%%%%%%%%%%%%%%%%%%%%%%%%%%%%%%
\subsubsection{Electrical Power Budget and Power Generation}
\label{sec:powerBudget}
%%%%%%%%%%%%%%%%%%%%%%%%%%%%%%%%%%%%%%%%%%%%%%%%%

\ac{JAO} currently operates with three turbines at the \ac{OSF}, each with a maximum output of 3.3~MW. Under normal conditions, one turbine is active while the remaining two are kept in standby. To maintain stability in the power generation system, the average power demand must be below 3~MW for the active turbine.
As of now, the current monthly average power demand, excluding the \ac{WSU} systems, ranges between 2.4~MW and 2.5~MW.   

During the \ac{WSU} parallel deployment phase, additional power will be required. The estimated increase includes 879~kW for equipment installed at \ac{OSF} \cite{OCRO_DSN}, and approximately 0.55~kW per antenna being retrofitted for \ac{WSU} (Table~\ref{tab:powerBudget}). This will add to the total of nominal 915~kW additional power required once all 66 antennas are retrofitted. 

%%%%%%%%%%%%%%%%
\begin{table}[p]
    \centering
    \caption{WSU Additional Electrical Budget}
    \label{tab:powerBudget}
    \begin{tblr}{
        colspec = {Q[l,4cm] Q[c,2cm] Q[c,2cm] Q[l,7cm]},
        cells={font=\footnotesize},
        row{1}={headercolor},
        row{2,10}={highlight2},
        row{9,18}={light_green},
        hlines,
        vlines
    }
    \textbf{Component} & \textbf{Average demand (kW)} & \textbf{UPS needed?} & \textbf{Notes/References} \\
    \SetCell[c=4]{c} AOS \\
    Front-End additional instruments & 0.1 & No & Additional power consumption will occur only with the new Band 2 receiver. \\
    \ac{WIFP}-\ac{IFSC} & 0.0 & No & The same power supply voltages from the \ac{FEPS} used for the existing IF switch are reused \cite{WIFPDescription}.\\
    \ac{WIFP}-\ac{DGOT} & 0.11 & No & \cite{WIFPDescription} \\
    \ac{WIFP}-\ac{DSPOT} & 0.34 & No & \cite{WIFPDescription} \\
    \ac{DTS}-\ac{AOS} Patch Panels & 0.0 & No & Additional EDFAs will be installed for future antennas located more than 15 km from AOS-TB \cite{dts_pdr_dsn}. \\
    Total additional power per antenna & 0.55 &  &  \\
    \textbf{Total additional power at AOS for 66 antennas} & \textbf{36.4} &  & The worst-case scenario assumes 66 antennas, each configured for parallel deployment. \\
    \SetCell[c=4]{c} OSF \\
    \ac{ATAC} & 332 & Yes & \cite{OCRO_DSN} \\
    \ac{DTS} at \ac{OCRO} & 10 & Yes & \cite{OCRO_DSN}, estimated as 7.7~kW in \cite{dts_pdr_dsn} \\
    \ac{TPGS} & 40 & Yes & \cite{OCRO_DSN} \cite{TPGS_cosdd} \\
    Computing infrastructure & 44 & Yes & \cite{OCRO_DSN} \\
    OCRO auxiliaries & 30.5 & Yes & \cite{OCRO_DSN} \\
    UPS room loads & 70 &  & \cite{OCRO_DSN} \\
    HVAC for OCRO & 352 &  & \cite{OCRO_DSN} \\
    \textbf{Total additional power at OSF} & \textbf{879} &  & \cite{OCRO_DSN} \\
    \end{tblr}
\end{table}
%%%%%%%%%%%%%%%%

At the end of 2024, a load test was conducted, and it was demonstrated that each turbine is capable of sustaining a stable operation up to 3.7~MW of load. Taking into account the current average load of 2.5~MW, and the nominal additional power of 0.915~MW required during the parallel deployment phase, a single turbine should be sufficient to sustain a stable operation running at 92\% of its empirical capacity (see \cite{WSU_rollout_plan} for more details). If this additional demand causes the total load to exceed the stable operating margin of a single turbine, mitigation strategies will be required to avoid activating a second turbine. Operating a second turbine for only a small increase in load is considered inefficient and is not the preferred solution. 

The 915~kW additional power is the nominal worst-case scenario, in which the primary and secondary chillers of OCRO are running at full capacity, OCRO's UPS is empty and charging, and all the 66 retrofitted antennas, TPGS, and computing infrastructure are in full operation. A more realistic estimation was presented in
Fig.~\ref{fig:WSU_incremental_power_breakdown_analysis}, which shows the progressive increase in the total additional power required by the WSU during the \ac{AIVC} and \ac{TWO} periods, as the number of retrofitted antennas integrated into the array increases. In this scenario, a maximum of the total additional power required is $\sim$650kW, which represents a 85\% of the empirical capacity of a single turbine. This figure also shows sporadic power demands during 2027 for the verification of the OCRO room in preparation for its acceptance (green spikes of $\sim$200~kW). 


%%%%%%%%%%%%%%%%
\begin{figure}[t!]
    \centering
    \includegraphics[width=\textwidth]{Figures/WSU-incremental-power-breakdown-analysis.png}
    \caption{WSU incremental power breakdown analysis.}
    \label{fig:WSU_incremental_power_breakdown_analysis}
\end{figure}
%%%%%%%%%%%%%%%%

%%%%%%%%%%%%%%%%
\begin{figure}[hbt!]
    \centering
    \includegraphics[width=\textwidth]{Figures/WSU-incremental-power-drawn.png}
    \caption{Additional WSU power consumption during parallel deployment and after the stop of legacy operations after Milestone 1. This chart assumes that 4$\times$ correlator capabilities will be deployed from late 2028 and will have finalized deployment by October 2030.}
    \label{fig:WSU_power_consumption_in_milestone_M1}
\end{figure}
%%%%%%%%%%%%%%%%

Once the \ac{WSU} deployment is complete, power demand estimates will take into account (i) the increased consumption from new \ac{WSU} subsystems and (ii) the reduced consumption resulting from retiring legacy components. As shown in Fig.~\ref{fig:WSU_power_consumption_in_milestone_M1}, once WSU-enabled scientific operations have started, the \acp{LRU} that comprise the legacy signal chain will be turned off. In an antenna, a 0.63~kW of power reduction will be obtained when retiring the legacy \ac{IFP}, \ac{LO2}, \ac{DGCK}, \ac{DTX}, and the \ac{ABM}. 



% reduced consumption: see TWO report page 78
% - IFP, LO2, DGCK, DTX, and ABM per AE: 2*83 + 24.4*4 + 10.1 + 53.4*4 + 150 = 637.3W   
% then, adding power consumption of E2C: 6.8W$ 
% => per AE de reduced comsumption will be 637.3 - 6.8W = 630W  => 0.63kW 
% 0.63 * 66 = 41.58 kW ~ 42kW 
% - BLC, ACAC, ACACSPEC: 264.7 kW in total (max load, from ICD)

In the technical building at the AOS, a 200~kW of power reduction will be obtained by shutting down the \ac{BLC} and the \ac{ACASPEC}.% (the \ac{ACAC} was already decommissioned before the start of the WSU AIVC process). 
In summary, a reduction of 42~kW (considering 66 antennas) and 200~kW can be obtained by shutting down the legacy equipment that is no longer needed, resulting in a total reduction of 242~kW. Therefore, \textbf{on average, the total additional power required to operate in the WSU era compared to current operation will be 401~kW and total power requirements can be met with only one turbine.} 






%%%%%%%%%%%%%%%%%%%%%%%%%%%%%%%%%%%%%%%%%%%%%%%%%
\subsubsection{OSF Computer Room Services}
\label{sec:osf_computer_room}
%%%%%%%%%%%%%%%%%%%%%%%%%%%%%%%%%%%%%%%%%%%%%%%%%
The OSF Computer Room is a 
%$\approx$ 200\,m$^2$  
datacenter located at the \ac{OSF}.  Its primary function is to host the following services:
\begin{itemize}[itemsep=-0.5em]
    \item Servers and computers that provide the offline and online infrastructure (Sections~\ref{sec:offline} and \ref{sec:online}, respectively).
    \item Network equipment required for OSF operations (Section~\ref{sec:opsInfrastructure}).
    \item The OSF data buffer system (Section~\ref{osf_data_buffer}).
    \item The OSF-to-SCO data transmission service (Section~\ref{ssub:osf_to_sco_data_transmission_and_buffering}), and
    \item Fiber interconnections at the OSF.
\end{itemize}

% Beside the computing and networking services provided by the OSF Computer Room, an additional key functionality to be hosted in it aims to interconnect all the subsystems of the telescope distributed in many places of the observatory to emerge the functionality of a radio telescope. The connections are done through fiber optic links. Based on the ALMA WSU design, both \ac{ATAC} and \ac{TPGS} will be located at the OSF, and housed inside of \ac{OCRO} (\ref{sec:OCRO}), sky  signals digitized by the signal chain of each antenna located at the AOS will be transported by the AOS-OSF \ac{FOS} (\ref{sec:FOS}) and terminated at the \ac{OCRO}, observing data acquired and processed by ATAC, TPGS must be sent and persisted in the OSF data buffer system (\ref{osf_data_buffer}). In addition, the current AOS-OSF fibers, OSF antenna pads fibers, and the legacy \ac{HILSE} fibers are terminated in the OSF Computer Room. All these fibers must be interconnected to establish the different scenarios needed during the AIVC, TWO and WSU operation phases.




% The OSF Computer Room will provide the fiber connectivity functionalities that used to be handled by the Patch Panel Room located at the \ac{AOS}. Fig.~\ref{fig:OSF_Computer_Room_and_OCRO} shows how the fibers will be interconnected between the different facilities at the Observatory. The Cable-05 and Cable-06 will provide the interconnection between the \ac{OCRO} and the OSF Computer Room. The Cable-05 will transport DTS signals originated by antennas located at the OSF pads to ATAC and TPGS, while Cable-06 will transport M\&C and data produced by ATAC and TPGS to the OSF data buffer system ((\ref{osf_data_buffer})). A fiber patch panel infrastructure will be installed in the OSF Computer Room to allow manually interconnect the corresponding fiber threads to support to the following possible use cases needed during the AIVC, TWO and WSU operation phases:

In addition to hosting computing and networking services, the OSF Computer Room plays a critical role in interconnecting the many distributed telescope subsystems. Both  \ac{ATAC} and \ac{TPGS}  will be located at the \ac{OSF} and housed inside the \ac{OCRO} (Section~\ref{sec:OCRO}). Sky signals digitized by each antenna at the \ac{AOS} will be transported via the \ac{FOS} (Section~\ref{sec:FOS}) and terminated at the \ac{OCRO}. Observing data acquired and processed by ATAC and TPGS will be sent to and stored in the OSF data buffer system (Section~\ref{osf_data_buffer}).

The OSF Computer Room will continue to serve as the termination point for several fiber systems, including the AOS-OSF fiber links, the fibers from the OSF antenna pads, and the legacy \ac{HILSE} fibers. These fiber connections must be flexibly interconnected to support a range of operational scenarios during the AIVC, TWO, and WSU phases.

Part of the fiber connectivity functions previously managed by the Patch Panel Room at the \ac{AOS} will now be handled by the OSF Computer Room. Figure~\ref{fig:OSF_Computer_Room_and_OCRO} illustrates the fiber interconnections between the OSF Computer Room and the OCRO. Two main cables provide this connectivity:
\begin{itemize}[itemsep=-0.5em]
    \item Cable-05, which carries DTS signals from antennas to ATAC and TPGS. Its fiber threads can be connected to Cable-03 to interconnect DTS signals of antennas located at the AOS or be connected to the Cable-02 to interconnect DTS signals of antennas located at the OSF.
    \item Cable-06, which transports monitor and control signals and science data from ATAC and TPGS to the OSF data buffer system (Section~\ref{osf_data_buffer}).
\end{itemize}

A dedicated fiber patch panel infrastructure will be installed in the OSF Computer Room to allow manual configuration of fiber connected as needed for the following use cases:
%interconnection of fibers, supporting the following use cases across the AIVC, TWO, and WSU operational phases:


\begin{itemize} [itemsep=-0.5em]
    \item Use Case 1: Legacy science operations.
    \item Use Case 2: WSU operations.
    \item Use Case 3: Legacy \ac{HILSE} operation.
    \item Use Case 4: Antenna maintenance at the OSF (during Legacy and WSU operation).
    \item Use Case 5: WSU AIVC Stage-2 and Stage-3. 
    \item Use Case 6: WSU AIVC Stage-4, and
    \item Use Case 7: WSU AIVC Stage-5. 
\end{itemize}


%%%%%%%%%%%%%%%%
\begin{figure}[hbt!]
    \centering
    \includegraphics[width=\textwidth]{Figures/OSF_computer_room_and_OCRO.png}
    \caption{Fiber interconnections between the OSF Computer Room and the \ac{OCRO}.}
    \label{fig:OSF_Computer_Room_and_OCRO}
\end{figure}
%%%%%%%%%%%%%%%%


%%%%%%%%%%%%%%%%%%%%%%%%%%%%%%%%%%%%%%%%%%%%%%%%%
\subsubsection{Legacy Fiber Infrastructure between AOS and OSF}
\label{sec:legacy_FOS}
%%%%%%%%%%%%%%%%%%%%%%%%%%%%%%%%%%%%%%%%%%%%%%%%%

The existing fiber optic cable between the \ac{AOS} and \ac{OSF} contains 48~fibers, organized into four 12-fiber mini-tubes \cite{aos_osf_fo_status}. Three mini-tubes connect the \ac{AOS} to the OSF Computing Room (Section~\ref{sec:osf_computer_room}), while the other two connect AOS-\ac{PPS} to OSF-\ac{PPS}.

Of the 36~fibers connecting AOS-TB and OSF-\acs{TF}, 21 remain available; the rest are currently used for IT communication infrastructure and for operating the current \ac{HILSE} subsystem. 
Along the AOS-PPS to OSF-PPS path, 6~fibers remain available, resulting in a total of 27 available fibers across all mini-tubes.

The current AOS–OSF fiber infrastructure will continue to support these essential services during WSU operations. These existing fibers will remain in place to connect the \ac{PPS} subsystem and other service links between the \ac{OSF} and \ac{AOS}. 
The \ac{FOS} installed for the WSU (Section~\ref{sec:FOS}) will be used exclusively for \ac{WSU} signal chain data transmission.
